\clearpage\thispagestyle{empty}
\pdfbookmark[0]{Quellen zur Bedeckung der Plejaden (29. August 1820)}{pleiades1824-sources}
\begin{center}\textbf{Quellen zur Bedeckung der Plejaden (29. August 1820)}\end{center}
\noindent Diese KI-Übertragung ins heutige Deutsch folgt dem Bericht \textit{Sternbedeckung} in \href{https://gdz.sub.uni-goettingen.de/id/PPN236007467}{den \textit{Werken}, Band~VI} (Göttingen, 1874), S.~448. Gauß teilte die Göttinger Beobachtungen in einem Brief mit. Die Werkausgabe nennt Bessel in eckigen Klammern als Empfänger. Die entsprechende Stelle steht in den \textit{Astronomischen Nachrichten}, Band~III, Nr.~50 (Februar 1824), \href{https://api.digitale-sammlungen.de/iiif/image/v2/bsb10049428_00023/full/1800,/0/default.jpg}{Sp.~22}, in der von Bessel mitgeteilten Untersuchung \textit{Berechnung verschiedener Sternbedeckungen} von Rosenberger, Strehlke und Klupsz. Die \href{https://api.digitale-sammlungen.de/iiif/image/v2/bsb10049428_00021/full/1800,/0/default.jpg}{Überschrift der Untersuchung} und die Textstelle stammen aus dem Exemplar der Bayerischen Staatsbibliothek. H.~C. Schumacher gab den Band 1825 in Altona heraus.
\par\smallskip\noindent Die Tabelle der Werkausgabe hat vier Zeilen, enthält zwei Wiederholungsstriche und setzt die Einheitenzeichen in die erste Zeile. Bei der ersten Zeitangabe für $\eta$ steht in der Zeitschrift an der Sekundenstelle 3,37 und in der Werkausgabe 37.37. Die Übertragung behält die Zahlen der Werkausgabe und die eingeklammerte Erläuterung zu $x$ bei. KI-Übertragung und KI-Satz.
\clearpage
